3d gaussian splatting (3DGS) is a relatively new method for recreating scene geometry represented by gaussian blobs inside of regular meshes learned from images of a scene from multiple viewpoints. Note that the problem it tries to solve differs from physics based rendering (PBR)methods. 3DGS is given images as input to recreate the scene to be finally rendered. PBR methods however are given the scene geometry first and tries to render photo-realistic image with it. The rendering part of 3DGS is built as a prerequisite to its back propagation steps: scene is rendered into image in order to be compared with its ground truth.This rendering pipeline uses splatting operation as well which is comparable to our method. 

The gaussian blobs are 3D gaussian spheres (ellipsoids) which their shapes are determined by mean and covariance (the two variables that can squeeze or strech these blobs) and the final contribution to the image is determined by their base color and alpha value. Then pipeline then uses alpha blending to sum up all the blobs for each pixel. The key differences is here the alphablending type is over that means the final color is composited by different blobs (show equation) and this operation iss not commutable. However in our method the blending operation is simply an addition where no alpha value is involved.
